% Generated draft: requires mathematical and source review.
\documentclass[11pt]{article}
\usepackage[margin=25mm]{geometry}
\usepackage{fontspec}
\setmainfont{Latin Modern Roman}
\usepackage{amsmath,amssymb,mathtools,unicode-math}
\setmathfont{Latin Modern Math}
\usepackage{enumitem,needspace,xurl,hyperref}
\setlength{\parindent}{0pt}
\setlength{\parskip}{5pt}
\setlength{\emergencystretch}{3em}
\newcommand{\N}{\mathbb{N}}
\newcommand{\Z}{\mathbb{Z}}
\newcommand{\Q}{\mathbb{Q}}
\newcommand{\R}{\mathbb{R}}
\newcommand{\C}{\mathbb{C}}
\newcommand{\F}{\mathbb{F}}
\newcommand{\A}{\mathbb{A}}
\newcommand{\PP}{\mathbb{P}}
\newcommand{\E}{\mathbb{E}}
\newcommand{\eps}{\varepsilon}
\newcommand{\dd}{\,\mathrm{d}}
\begin{document}
\section*{UnsolvedMath ID: 6700066}
\textbf{Source problem code:} catalogue entry
\subsection{Definitions and mathematical statement}
Let $Y\subset\R^n$ be a domain with sufficiently regular boundary, let $\varepsilon>0$, and let $k\in\{2,\ldots,n-1\}$. The main conjecture assumes
\[
 H_{\partial Y}\ge n-k+\varepsilon
\]
and asks for a continuous surjection $\Delta:Y_{-1}\to P^{k-1}$ onto a $(k-1)$-dimensional polyhedral space such that
\[
 \operatorname{diam}(\Delta^{-1}(p))\le C(n,\varepsilon)
 \qquad(p\in P^{k-1}).
\]
The bound is uniform over all such domains and fibers. The source does not fix the sign and sum-versus-average normalization of mean curvature, the meaning of the unit interior $Y_{-1}$, or whether fiber diameters use ambient Euclidean distance or an intrinsic metric.

For a metric space $M$, the supplied definition of macroscopic dimension is the least dimension of a polyhedral space admitting a continuous map from $M$ whose point fibers have uniformly bounded diameter, with some bound $d(M)$. Thus the asserted bound implies $\operatorname{dim}_{\mathrm{macro}}Y_{-1}\le k-1$ in that convention, with the stronger specified uniform constant. Properness or uniform control on inverse images of larger sets is not added.

The source proposes $Y=B^{n-k}(1)\times\R^k$ as the extremal $\varepsilon=0$ example and invokes Lebesgue's lemma to say that $\R^k$ cannot map continuously to a $(k-1)$-polyhedron with uniformly bounded point fibers. This example needs checking: in the usual summed-principal-curvature convention its cylindrical boundary has mean curvature $n-k-1$, not $n-k$; moreover a strict distance-greater-than-one interior of the unit-radius cylinder is empty, whereas a non-strict convention can retain its axis. No changed dimension or radius is inferred.

The passage additionally asserts that a connected domain with $H_{\partial Y}>n-2+\varepsilon$ has bounded $Y_{-1}$. This conclusion is retained separately, with the same missing conventions; its derivation is not supplied.

The extraction then begins a different section of definitions for an upcoming scalar-curvature analogue. For maps $\varphi:X\to S$ it gives
\[
 d_\infty(\varphi_1,\varphi_2)
       =\sup_{x\in X}d_S(\varphi_1(x),\varphi_2(x)),
\]
and equips subsets of the map space with the induced length metric, obtained by infimizing lengths of curves measured using $d_\infty$. It also describes this, except for unspecified pathological cases, through metrics locally agreeing with the sup metric. The length metric may exceed the sup metric when the subset is distorted.

For $C(X,S)$, the continuous-map space, the source claims that a locally isometric covering $\widetilde S\to S$ induces an ``isometry'' $C(X,\widetilde S)\to C(X,S)$. A literal global isometry is problematic because distinct deck-related lifts can have identical images; the intended component, quotient, or local formulation is not stated. Its footnote assumes the spaces locally contractible, such as manifolds or cell complexes, and $S$ a length space. It then claims for compact $S$ that $C(B^n,S)$ has infinite length-metric diameter if $\pi_1(S)$ is infinite and finite diameter if $\pi_1(S)$ is finite, and gives
\[
 \operatorname{diam}C(B^n,S^m(1))=\pi\quad(m>n),
 \qquad
 \operatorname{diam}C(B^n,S^n(1))\le3\pi.
\]
These are adjacent source claims with unspecified map-space conventions, not extra conclusions of the mean-curvature conjecture. The forthcoming scalar-curvature analogue itself is not included.
\subsection{Short English statement}
A positive margin above the stated boundary mean-curvature threshold should force the unit interior of a Euclidean domain to have bounded-fiber maps into a lower-dimensional polyhedron, with a bound depending only on dimension and the margin. The source also states a bounded-interior consequence and includes later map-space material.

The supplied summary reports partial mean-convex-domain/macroscopic-dimension and filling-radius context, but no exact results or references resolving this conjecture. Its partial-progress label is retained with current status marked for review. The cylinder normalization and the adjacent covering-map claim have formulation issues that are not silently repaired.
\subsection{Sources}
\begin{itemize}
\item Local UnsolvedMath registry, permanent ID 6700066, source code catalogue entry, scalar-curvature Question [?70], proposed by Misha Gromov, 2017. Dataset curation and metadata: Ulam AI and UnsolvedMath Contributors, CC BY 4.0.
\item Misha Gromov, \emph{101 Questions, Problems and Conjectures around Scalar Curvature} (2017), printed tag [?70], occurrence 69, PDF page 70, including the extracted following map-space discussion and footnote 65, \url{https://www.ihes.fr/~gromov/wp-content/uploads/2018/08/101-problemsOct1-2017.pdf}; normalized pdf-text extraction, accessed 2026-07-29.
\item Lebesgue's lemma is invoked in the cylinder discussion without a separate citation; no exact later mean-curvature/macroscopic-dimension theorem is supplied.
\item Source-specific provenance: the author-hosted PDF is incomplete and publicly available; current status and statement-reuse terms NEEDS\_REVIEW, separately from the dataset licence.
\end{itemize}
\end{document}
